\documentclass{beamer}
\usepackage[utf8]{inputenc}

\usetheme{Madrid}
\usecolortheme{default}
\usepackage{amsmath,amssymb,amsfonts,amsthm}
\usepackage{txfonts}
\usepackage{tkz-euclide}
\usepackage{listings}
\usepackage{adjustbox}
\usepackage{array}
\usepackage{tabularx}
\usepackage{gvv}
\usepackage{lmodern}
\usepackage{circuitikz}
\usepackage{tikz}
\usepackage{graphicx}
\usepackage{multicol}
\setbeamertemplate{page number in head/foot}[totalframenumber]

\usepackage{tcolorbox}
\tcbuselibrary{minted,breakable,xparse,skins}



\definecolor{bg}{gray}{0.95}
\DeclareTCBListing{mintedbox}{O{}m!O{}}{%
  breakable=true,
  listing engine=minted,
  listing only,
  minted language=#2,
  minted style=default,
  minted options={%
    linenos,
    gobble=0,
    breaklines=true,
    breakafter=,,
    fontsize=\small,
    numbersep=8pt,
    #1},
  boxsep=0pt,
  left skip=0pt,
  right skip=0pt,
  left=25pt,
  right=0pt,
  top=3pt,
  bottom=3pt,
  arc=5pt,
  leftrule=0pt,
  rightrule=0pt,
  bottomrule=2pt,

  colback=bg,
  colframe=orange!70,
  enhanced,
  overlay={%
    \begin{tcbclipinterior}
    \fill[orange!20!white] (frame.south west) rectangle ([xshift=20pt]frame.north west);
    \end{tcbclipinterior}},
  #3,
}
\lstset{
    language=C,
    basicstyle=\ttfamily\small,
    keywordstyle=\color{blue},
    stringstyle=\color{orange},
    commentstyle=\color{green!60!black},
    numbers=left,
    numberstyle=\tiny\color{gray},
    breaklines=true,
    showstringspaces=false,
}
%------------------------------------------------------------
%This block of code defines the information to appear in the
%Title page
\title %optional
{2.0.12}
%\subtitle{A short story}

\author % (optional)
{BEERAM MADHURI - EE25BTECH11012}

\begin{document}


\frame{\titlepage}
\begin{frame}{Question}
A product is made by mixing three raw materials 1, 2, 3 in varying proportions, where material 1 must account for not more than 50\% of the total. If $x, y$ and $z$ are the amounts of raw materials 1, 2, 3, respectively, this constraint can be modeled as \hfill(2007 PI)

\begin{enumerate}[label=\alph*)]
    \item $x \le 0.5$
    \item $x \le 0.5(x + y + z)$
    \item $0.5x \le x + y + z$
    \item $x \ge 0.5(y + z)$
\end{enumerate}
\end{frame}

\begin{frame}{solution}
Given,\\
Raw Materials: 1,2,3\\
Amount: x,y,z\\
let,\\
Total amount of Product: T\\
\begin{align}
    T = x+y+z
\end{align}
\end{frame}
\begin{frame}
given constraint: Material one (x) must not account for more than 50\% of the total.\\
\begin{align} 
    x \le 0.5(T)\\
    x \le 0.5(x+y+z)\\
    x\le y+z
\end{align}

Correct Option: b
\end{frame}

\begin{frame}[fragile]
    \frametitle{Python Code}
    \begin{lstlisting}
import numpy as np
import matplotlib.pyplot as plt
from mpl_toolkits.mplot3d import Axes3D

y = np.linspace(0, 10, 20)
z = np.linspace(0, 10, 20)
Y, Z = np.meshgrid(y, z)

# The boundary condition: x = y + z (which is x = 0.5(x + y + z))
X_boundary = Y + Z

# Setup the figure
fig = plt.figure(figsize=(10, 8))
ax = fig.add_subplot(111, projection='3d')
\end{lstlisting}
\end{frame}

\begin{frame}[fragile]
\frametitle{Python Code}
\begin{lstlisting}
# 1. Plot the Boundary Plane
surf = ax.plot_surface(X_boundary, Y, Z, alpha=0.5, cmap='viridis', edgecolor='none')

# 2. Visualize the Feasible Region (Constraint)
# We generate random points and keep only those where x <= y + z
x_pts = np.random.uniform(0, 20, 5000)
y_pts = np.random.uniform(0, 10, 5000)
z_pts = np.random.uniform(0, 10, 5000)

# Filter points that satisfy the constraint: x <= 0.5(x + y + z)
mask = x_pts <= (y_pts + z_pts)
x = np.linspace(0, 20, 30)
y = np.linspace(0, 10, 30)
z = np.linspace(0, 10, 30)
\end{lstlisting}
\end{frame}

\begin{frame}[fragile]
\frametitle{Python Code}
\begin{lstlisting}
X, Y, Z = np.meshgrid(x, y, z)

mask = X <= (Y + Z)
ax.scatter(X[mask], Y[mask], Z[mask],
           color='green', s=2, alpha=0.2)

# Labels and Title
ax.set_xlabel('Amount of X (Material 1)')
ax.set_ylabel('Amount of Y (Material 2)')
ax.set_zlabel('Amount of Z (Material 3)')
ax.set_title('Constraint Region: x $\leq$ 0.5(x + y + z)')

plt.show()
\end{lstlisting}
\end{frame}

\begin{frame}{Plot}
\begin{figure}[H]
\centering
\includegraphics[width=0.75\columnwidth]{Figures/plot.png}
\caption{Boundary Condition}
\label{fig:placeholder}
\end{figure}
\end{frame}

\end{document}